% =====================================================================
% Clase -- Cálculo 11° -- Trimestre I -- Semana 03
% Sesiones 006 (lun 14 sep., 100 min), 007 (jue 17 sep., 55 min), 008 (vie 18 sep., 50 min)
% Tema 2 de la guía: Los irracionales y la recta real
% Material del docente (no se entrega a los estudiantes).
% =====================================================================
\documentclass[11pt]{article}
\makeatletter\def\input@path{{../../../../../plantillas/estilo/}}\makeatother
\usepackage{mmcantillo}
\guiadelestudiante{../../guia-didactica/guia-periodo-I-numeros-reales}

\tipodocumento{Clase}
\asignatura{Cálculo}
\titulo{Semana 3: $\sqrt{2}$, los irracionales notables y la recta real}
\grado{11°}
\periodo{I}
% DBA: matematicas grado 11 · DBA 1 — Utiliza las propiedades de los números (naturales,
%      enteros, racionales y reales) y sus relaciones y operaciones para construir y comparar
%      los distintos sistemas numéricos.
% Evidencias trabajadas: (1) describe propiedades de los números y las operaciones que son
%      comunes y diferentes en los distintos sistemas numéricos (¿es cerrada la suma de
%      irracionales?); (3) construye representaciones de los conjuntos numéricos y establece
%      relaciones acordes con sus propiedades ($\sqrt{2}$ en la recta, $\R = \Q \cup \mathbb{I}$).
% Nota: el plan trae «Los sistemas numéricos» (semanas 01-02) sin dictar; dónde se recupera
%      está pendiente de decisión de la docente. Esta semana sigue el plan tal como está.

\begin{document}

\encabezado*

\begin{center}
{\renewcommand{\arraystretch}{1.3}
\begin{tabularx}{\linewidth}{@{}l l l X l@{}}
  \toprule
  \textbf{Sesión} & \textbf{Fecha} & \textbf{Min} & \textbf{Subtema} & \textbf{Entregables}\\
  \midrule
  006 & lun 14 sep., 07:50 & 100 & $\sqrt{2}$ no es racional: la crisis de Hipaso & Taller 3\\
  007 & jue 17 sep., 06:55 & 55  & Irracionales notables: $\pi$, $e$, $\sqrt{n}$ & —\\
  008 & vie 18 sep., 13:40 & 50  & La recta real (Dedekind) & Tarea 3 (para el lun 21 sep.)\\
  \bottomrule
\end{tabularx}}
\end{center}

Referencia: \refguia{tema:irracionales}.

% ---------------------------------------------------------------------
\seccion{Sesión 006 — lunes 14 de septiembre · 07:50 · 100 min}

\textbf{Objetivo.} Demostrar por contradicción que $\sqrt{2}$ no es un número racional y
reconocer que la recta tiene puntos que ninguna fracción alcanza.

\begin{center}
{\renewcommand{\arraystretch}{1.3}
\begin{tabularx}{\linewidth}{@{}l l X@{}}
  \toprule
  \textbf{Momento} & \textbf{Min} & \textbf{Actividad}\\
  \midrule
  Situación & 0–15 & Cada estudiante dibuja un cuadrado de lado 1 (un cuadro de 10
    cuadritos del cuaderno) y mide su diagonal con regla: $\approx 1{,}4$. Pregunta: ¿qué
    fracción es \emph{exactamente}? Prueban con $\frac{7}{5}$ y $\frac{3}{2}$ elevando al
    cuadrado: $1{,}96$ y $2{,}25$; ninguna da 2.\\
  Historia & 15–25 & El hilo de la guía: los pitagóricos creían que todo es razón de
    enteros; según la leyenda, Hipaso de Metaponto reveló que la diagonal no lo es. (Es una
    leyenda: las fuentes antiguas no coinciden en los detalles.)\\
  Demostración & 25–50 & En el tablero, la demostración de la guía paso a paso. Antes de cada
    paso pregunta al grupo: ¿por qué si $p^2$ es par, $p$ es par? (Si $p$ fuera impar,
    $p = 2m+1$ y $p^2 = 4m^2 + 4m + 1$ sería impar.) ¿Dónde está la contradicción?\\
  Taller 3 & 50–90 & En parejas.\\
  Cierre & 90–100 & $\R = \Q \cup \mathbb{I}$: los decimales infinitos no periódicos. Pregunta
    para el jueves: ¿conocen otro número que no sea una fracción?\\
  \bottomrule
\end{tabularx}}
\end{center}

\begin{nota}
Para ver la estructura de la demostración, escribe en una columna aparte lo que se
\emph{supone} (la fracción es irreducible) y lo que se \emph{deduce} ($p$ y $q$ pares): la
contradicción es que las dos columnas no pueden ser ciertas a la vez. El ejercicio de la guía
«Adapta la demostración para $\sqrt{3}$» queda como reto voluntario.
\end{nota}

% ---------------------------------------------------------------------
\seccion{Sesión 007 — jueves 17 de septiembre · 06:55 · 55 min}

\textbf{Objetivo.} Reconocer los irracionales más usados ($\pi$, $e$, las raíces no exactas
$\sqrt{n}$) y distinguir un valor exacto de sus aproximaciones racionales.

\begin{center}
{\renewcommand{\arraystretch}{1.3}
\begin{tabularx}{\linewidth}{@{}l l X@{}}
  \toprule
  \textbf{Momento} & \textbf{Min} & \textbf{Actividad}\\
  \midrule
  Revisión & 0–10 & Se devuelve el Taller 3 y se comenta el punto 5: la suma de dos
    irracionales puede ser racional, así que $\mathbb{I}$ no es cerrado para la suma
    (a diferencia de $\Q$).\\
  $\sqrt{n}$ & 10–20 & Si $n$ no es un cuadrado perfecto, $\sqrt{n}$ es irracional (se acepta
    sin demostrar; el caso $\sqrt{3}$ sale igual que $\sqrt{2}$). Clasificar en el tablero
    $\sqrt{16}$, $\sqrt{17}$, $\sqrt{0{,}25}$, $\sqrt{8}$.\\
  $\pi$ & 20–35 & $\pi$ es la razón entre la longitud de una circunferencia y su diámetro.
    Lambert demostró en 1761 que es irracional. Aproximaciones: $3$ (la Biblia), $\frac{22}{7}$
    (Arquímedes), $\frac{355}{113}$ (Zu Chongzhi). Si hay tiempo, medir con una cuerda el
    contorno y el diámetro de un objeto circular del salón.\\
  $e$ & 35–48 & $e \approx 2{,}71828$: en la calculadora, $\left(1 + \frac{1}{n}\right)^n$ con
    $n = 1, 10, 100, 1000, 10\,000$ (dan $2$; $2{,}594$; $2{,}705$; $2{,}717$; $2{,}718$).
    Es el interés compuesto «cada instante»; volverá en el trimestre II con la función
    exponencial. Euler usó la letra $e$ y demostró su irracionalidad (1737).\\
  Cierre & 48–55 & Diferencia entre «$\pi = 3{,}14$» (falso) y «$\pi \approx 3{,}14$»
    (verdadero). Cada estudiante escribe en su cuaderno un racional y un irracional entre 3 y
    4.\\
  \bottomrule
\end{tabularx}}
\end{center}

% ---------------------------------------------------------------------
\seccion{Sesión 008 — viernes 18 de septiembre · 13:40 · 50 min}

\textbf{Objetivo.} Entender que cada número real es un punto de la recta y cada punto de la
recta es un número real (la recta «no tiene huecos»), y ubicar irracionales en ella. Es la
última hora del viernes: sesión activa.

\begin{center}
{\renewcommand{\arraystretch}{1.3}
\begin{tabularx}{\linewidth}{@{}l l X@{}}
  \toprule
  \textbf{Momento} & \textbf{Min} & \textbf{Actividad}\\
  \midrule
  Construcción & 0–20 & Con regla y compás, cada estudiante lleva $\sqrt{2}$ a la recta (la
    figura de la guía) y, con un triángulo de catetos 1 y 2, también $\sqrt{5}$. Luego ubica
    $1 + \sqrt{2}$ trasladando el segmento.\\
  Dedekind & 20–35 & Idea de la cortadura (Dedekind, \emph{Continuidad y números
    irracionales}, 1872): partir los racionales en dos grupos, $A = \{$racionales negativos o
    con cuadrado menor que 2$\}$ y $B$ el resto. $A$ no tiene máximo y $B$ no tiene mínimo:
    en el corte hay un «hueco» de $\Q$, y ese hueco es $\sqrt{2}$. Al llenar todos los huecos
    se obtiene $\R$.\\
  Juego & 35–45 & «¿Hay hueco?»: la docente dice dos números ($1{,}41$ y $1{,}42$; $\frac{1}{3}$
    y $0{,}334$) y el grupo que primero dé un número entre ellos gana un punto. Prepara la
    densidad del lunes.\\
  Cierre & 45–50 & Se asigna la Tarea 3 para el lunes 21.\\
  \bottomrule
\end{tabularx}}
\end{center}

\begin{nota}
Lleva compás y regla para el tablero, y pide a los estudiantes el suyo el jueves. La cortadura
se presenta como idea, sin formalizar: basta con que vean que $A$ y $B$ no tienen «borde»
racional.
\end{nota}

\end{document}
